% Book 3, Year 1, chapter 8 -- the shell.
% Provenance (never printed): sources/book3/BRIEFS.md, chapter 8.
% ledger: b1:pipes-1972, b1:shell-user-program, b1:bourne-bash, b1:posix-expansion-order
\chapter{The Shell}\label{ch:b1:the-shell}

How many loans did each member of the town library make in 2025? A Python program of fifteen lines answers it; so does
one line typed at a terminal, assembled from five small programs that each do one thing. This chapter is about that
second way of working: the \omterm{def:b1:the-shell:shell}{shell}, the programs it starts, the streams that join them, and the scripts that turn a
line that worked once into a tool that works every time. The C programs of the previous chapters are among the
programs it starts.

\begin{recall}
The operating system's kernel runs programs and keeps their files in a tree of directories, each file named by a
path, absolute from the root \code{/} or relative to the current directory. A running program is a process. Each file
has permissions saying who may read, write or execute it.
\end{recall}

\section{Commands, processes and exit status}\label{sec:b1:the-shell:commands}

\begin{definition}[Shell]\label{def:b1:the-shell:shell}
A \emph{shell}\index{shell} is a program that reads command lines, typed or from a file, and runs them: for each
command it starts the named program with the given arguments, waits for it to end, and reads the next line. The shell
of this chapter is \code{bash}.
\end{definition}

\begin{history}[The shell and the pipe]
On the first Unix system, the command \omterm{def:b1:compilation-pipeline:interpreter}{interpreter} was not part of the operating system but an ordinary program
stored in a file, which any user could replace: \omterm{def:b1:the-shell:redirection}{redirection} with \code{<} and \code{>} came almost at once. Pipes came
in 1972, at the insistence of Douglas McIlroy, who had long argued that programs should be joined like operators.
The \omterm{def:b1:the-shell:shell}{shell} written by Stephen Bourne for the Seventh Edition of Unix fixed the language that \code{bash}, the GNU
project's ``Bourne-Again SHell'', still speaks.
% ledger: b1:vt100
Its users typed at terminals: a keyboard and a screen of text joined to the computer by a serial line. The VT100 of
1978 was one of the first to obey the standard escape sequences for moving the cursor, and the terminal windows of
today still obey them.
\end{history}

\omimg{book3/vt100-terminal.jpg}{0.42\linewidth}{A VT100 terminal, made by Digital Equipment Corporation: a keyboard
and a screen of text, with no computer of its own to run programs on. Photograph: Jason Scott, CC BY 2.0.}

\begin{definition}[Command-line argument]\label{def:b1:the-shell:argument}
The \omterm{def:b1:the-shell:shell}{shell} cuts a command line into words. The first names the program; the others are its \emph{command-line
arguments}\index{command-line argument}, which the program receives as strings: in C, as \code{argv[1]} to
\code{argv[argc - 1]}, \code{argv[0]} being the name it was started under.
\end{definition}

\begin{definition}[Exit status]\label{def:b1:the-shell:exit-status}
The \emph{exit status}\index{exit status} of a program is the small integer, from 0 to 255, it hands back when it ends:
the value returned by \code{main} or passed to \code{exit}. By convention 0 means success and any other value failure.
The \omterm{def:b1:the-shell:shell}{shell} keeps the status of the last command in the variable \code{$?}.
\end{definition}

\omoutput{code/bachelor-1/08-the-shell/outputs/status.txt}{Success is 0; \code{ls} reports a missing file with 2;
\code{grep -q} says only, by its status, whether it found a line.}\label{out:b1:the-shell:status}

\code{a && b} runs \code{b} only if \code{a} succeeded; \code{a || b} only if \code{a} failed; \code{a; b} runs both.
Programs are judged by their status, not by their messages: \code{grep -q} prints nothing at all.

\begin{definition}[Environment variable]\label{def:b1:the-shell:environment}
An \emph{environment variable}\index{environment variable} is a pair \code{NAME=value} in the environment of a running
program, a list of strings that it receives when it starts, as a copy of its parent's exported variables. A \omterm{def:b1:the-shell:shell}{shell}
variable becomes part of the environment of the programs the \omterm{def:b1:the-shell:shell}{shell} starts when it is \emph{exported}; \code{NAME=value
command} sets it for that one command. A C program reads it with \code{getenv("NAME")}.
\end{definition}

\omcode{code/bachelor-1/08-the-shell/listings/env_child.c}{2}{13}{A program that shows what it
received}\label{lst:b1:the-shell:env-child}

\omoutput{code/bachelor-1/08-the-shell/outputs/env.txt}{Quotes make \code{two words} one argument; a variable reaches
the program only once exported.}\label{out:b1:the-shell:env}

Every program documents its arguments: \code{man ls} shows the manual page of \code{ls}, \texttt{ls -{}-{}help} a summary.
A running program is stopped by Ctrl-C, which the terminal turns into a signal sent to it; signals are studied in the
Year~2 volume.

\section{Redirection and pipes}\label{sec:b1:the-shell:pipes}

Every program starts with three open streams (\cref{ch:b1:c-strings-structs-files}): standard input, standard output
and standard error. In a terminal all three are the terminal; the \omterm{def:b1:the-shell:shell}{shell} can connect them elsewhere.

\begin{definition}[Redirection]\label{def:b1:the-shell:redirection}
A \emph{redirection}\index{redirection} connects a standard stream of a command to a file: \code{> f} sends standard
output to \code{f}, emptied first; \code{>> f} appends to it; \code{< f} reads standard input from it; \code{2> f}
sends standard error; and \code{2>&1} sends standard error wherever standard output goes at that point of the line.
Redirections are performed by the \omterm{def:b1:the-shell:shell}{shell} before the program starts, from left to right.
\end{definition}

\omoutput{code/bachelor-1/08-the-shell/outputs/redirect.txt}{The error message ignores \code{>}; \code{2>} and
\code{2>&1} catch it.}\label{out:b1:the-shell:redirect}

\begin{omfigure}
\begin{tikzpicture}[p/.style={draw, rounded corners=2pt, minimum width=16mm, minimum height=8mm, font=\small\ttfamily},
                    f/.style={draw, minimum height=6mm, font=\scriptsize, inner sep=2pt}, >={Stealth[length=2mm]}]
\node[font=\small\bfseries] at (2,1.4) {\code{ls ... > all.txt}};
\node[p] (a) at (0.6,0) {ls};
\node[f] (t1) at (3.4,0.4) {all.txt};
\node[f] (t2) at (3.4,-0.4) {terminal};
\draw[->, thick] (a.east) ++(0,0.15) -- node[above, font=\scriptsize, pos=0.4] {1} (t1.west);
\draw[->, thick, omThm] (a.east) ++(0,-0.15) -- node[below, font=\scriptsize, pos=0.4] {2} (t2.west);
\node[font=\small\bfseries] at (8,1.4) {\code{ls ... > all.txt 2>&1}};
\node[p] (b) at (6.6,0) {ls};
\node[f] (u1) at (9.4,0) {all.txt};
\draw[->, thick] (b.east) ++(0,0.15) -- node[above, font=\scriptsize, pos=0.4] {1} (u1.west);
\draw[->, thick, omThm] (b.east) ++(0,-0.15) -- node[below, font=\scriptsize, pos=0.4] {2} (u1.west);
\end{tikzpicture}
\par\smallskip{\small Standard output (1) and standard error (2) of \code{ls}. With \code{> all.txt} alone, the error
still reaches the terminal; \code{2>&1}, written after it, sends the error where the output goes now, to the file.
Written before \code{> all.txt}, it would copy the old destination, the terminal (\cref{exo:b1:the-shell:6}).}
\end{omfigure}

\begin{definition}[Pipeline]\label{def:b1:the-shell:pipeline}
A \emph{pipeline}\index{pipeline (shell)} is a sequence of commands joined by the operator \code{|}: the \omterm{def:b1:the-shell:shell}{shell} starts
them all at once, and connects the standard output of each to the standard input of the next, so that data flows
through them without intermediate files. The \omterm{def:b1:the-shell:exit-status}{exit status} of a pipeline is that of its last command.
\end{definition}

\begin{definition}[Filter]\label{def:b1:the-shell:filter}
A \emph{filter}\index{filter (shell)} is a program that reads standard input, transforms it, and writes standard
output, so that it can stand anywhere in a \omterm{def:b1:the-shell:pipeline}{pipeline}: \code{sort} sorts lines, \code{uniq -c} collapses runs of equal
adjacent lines and counts them, \code{wc -l} counts lines, \code{head} and \code{tail} keep the first or last lines,
\code{cut} keeps fields, \code{tr} translates characters, \code{paste} joins lines.
\end{definition}

\omoutput{code/bachelor-1/08-the-shell/outputs/loans.txt}{The opening question: the five members with the most loans
in 2025, and the number of members who borrowed.}\label{out:b1:the-shell:loans}

\begin{omfigure}
\begin{tikzpicture}[p/.style={draw, rounded corners=2pt, minimum height=8mm, font=\scriptsize\ttfamily, inner sep=3pt},
                    >={Stealth[length=2mm]}, node distance=6mm]
\node[p] (a) {tail -n +2 loans.csv};
\node[p, right=of a] (b) {cut -d, -f3};
\node[p, right=of b] (c) {sort};
\node[p, right=of c] (d) {uniq -c};
\node[p, right=of d] (e) {sort -rn};
\node[p, right=of e] (f) {head -5};
\foreach \x/\y in {a/b, b/c, c/d, d/e, e/f} {\draw[->, thick] (\x) -- (\y);}
\draw[->, thick] (f.east) -- ++(0.6,0) node[right, font=\scriptsize] {terminal};
\foreach \x in {a,b,c,d,e,f} {\draw[->, omThm] (\x.south) -- ++(0,-0.55);}
\node[font=\scriptsize, omThm] at (5.2,-1.2) {each standard error still goes to the terminal};
\node[font=\scriptsize, align=center] at (1.0,0.9) {header dropped};
\node[font=\scriptsize, align=center] at (3.25,0.9) {member ids};
\node[font=\scriptsize, align=center] at (6.6,0.9) {\code{12 55}, \dots};
\end{tikzpicture}
\par\smallskip{\small The \omterm{def:b1:the-shell:pipeline}{pipeline} of \cref{out:b1:the-shell:loans}: six programs running at once, each reading what the
one before writes. \code{uniq -c} counts only adjacent equal lines, hence the \code{sort} before it.}
\end{omfigure}

A \omterm{def:b1:the-shell:pipeline}{pipeline} that answers a question deserves a check. The Python program below computes the same table from the same
file; the chapter's test runs both and compares them.

\omcode{code/bachelor-1/08-the-shell/listings/loans_check.py}{3}{20}{The same answer in
Python}\label{lst:b1:the-shell:loans-check}

\omoutput{code/bachelor-1/08-the-shell/outputs/loans_py.txt}{Identical, ties included: \code{sort -rn} orders equal
counts by the whole line, in reverse, which the key of \code{sorted} reproduces.}

\section{Patterns: globs and regular expressions}\label{sec:b1:the-shell:patterns}

Before running a command, the \omterm{def:b1:the-shell:shell}{shell} rewrites its words.

\begin{definition}[Pathname expansion, quoting, word splitting]\label{def:b1:the-shell:expansion}
\emph{Pathname expansion}\index{pathname expansion}\index{glob} replaces a word containing \code{*} (any string),
\code{?} (any one character) or \code{[...]} (one character of a set) by the sorted names of the files that match it, a
\emph{glob}; a pattern that matches nothing is left as it is, and names starting with a dot are matched only by a
pattern starting with a dot. \emph{Word splitting}\index{word splitting} cuts the result of an unquoted variable
expansion into several words at spaces, tabs and newlines. \emph{Quoting}\index{quoting} suspends both: between single
quotes every character is literal; between double quotes \code{$} still expands variables but no splitting or
pathname expansion happens; a backslash makes the next character literal.
\end{definition}

\omoutput{code/bachelor-1/08-the-shell/outputs/globs.txt}{\omterm{def:b1:the-shell:expansion}{Globs} expand to file names; quoted, they stay patterns;
braces are a \code{bash} expansion that makes words without looking at files.}\label{out:b1:the-shell:globs}

\omoutput{code/bachelor-1/08-the-shell/outputs/quoting.txt}{A file name with a space, unquoted and
quoted.}\label{out:b1:the-shell:quoting}

\begin{omfigure}
\begin{tikzpicture}[s/.style={draw, rounded corners=2pt, font=\scriptsize, align=center, minimum height=11mm,
                    text width=19mm, inner sep=2pt}, >={Stealth[length=2mm]}, node distance=4mm]
\node[s, fill=black!5] (b) {braces\\\code{{a,b}.csv}};
\node[s, fill=omCode!8, right=of b] (v) {variables, \code{$(...)}, \code{$((...))}};
\node[s, fill=omCode!8, right=of v] (w) {word splitting (unquoted results)};
\node[s, fill=omDef!12, right=of w] (g) {pathname expansion\\\code{*.csv}};
\node[s, fill=omThm!12, right=of g] (q) {quote removal};
\foreach \x/\y in {b/v, v/w, w/g, g/q} {\draw[->, thick] (\x) -- (\y);}
\node[font=\scriptsize, anchor=north west, align=left] at (b.south west) {\code{bash} only};
\end{tikzpicture}
\par\smallskip{\small The order in which the \omterm{def:b1:the-shell:shell}{shell} rewrites a command line (the tilde \texttt{\textasciitilde} is expanded
with the variables). Quotes act at every step, braces included, and are removed last: that is why \code{"$f"} reaches
the program as one word with the variable's value, spaces and all.}
\end{omfigure}

A pattern of a different kind describes the lines of a file. \code{grep} prints the lines that match a regular
expression; \code{sed} edits a stream line by line, here with its substitution \code{s/pattern/replacement/}; \code{find}
walks a tree of directories and lists the files that pass its tests. Regular expressions are used here as the
syntax of these tools: \code{.} is any character, \code{[0-9]} one digit, \code{x*} any number of \code{x}, \code{^}
and \code{$} the start and end of a line, \code{a|b} either, and with \code{-E} parentheses group and \code{{4}}
repeats. Their theory, and the automata that run them, belong to the Year~2 volume.

\omoutput{code/bachelor-1/08-the-shell/outputs/grep_sed.txt}{Counting, selecting and rewriting lines of the weather
file.}\label{out:b1:the-shell:grep-sed}

The two empty-or-error patterns, \code{,,} and \code{-99}, find the eight bad temperatures that
\cref{ch:b1:c-strings-structs-files} counted. In the last command, \texttt{\textbackslash{}3} is the text matched by the
third group: the date is rewritten day first.

\omoutput{code/bachelor-1/08-the-shell/outputs/find.txt}{\code{find} on a small tree: directories, files by name, by
size, and a command run on all of them.}\label{out:b1:the-shell:find}

\section{Shell scripts}\label{sec:b1:the-shell:scripts}

\begin{definition}[Shell script]\label{def:b1:the-shell:script}
A \emph{shell script}\index{shell script} is a file of \omterm{def:b1:the-shell:shell}{shell} commands run as a program. Its first line, \texttt{\#!/bin/bash},
names the \omterm{def:b1:compilation-pipeline:interpreter}{interpreter} that runs it; made executable, it is started like any program. Inside it, \code{$1}, \code{$2},
\dots{} are its arguments, \texttt{\$\#} their number and \code{"$@"} all of them, each kept as one word; its \omterm{def:b1:the-shell:exit-status}{exit status}
is that of \code{exit}, or of its last command.
\end{definition}

\omcode{code/bachelor-1/08-the-shell/listings/check_csv.sh}{1}{27}{A script that checks the town's
tables}\label{lst:b1:the-shell:check-csv}

\omoutput{code/bachelor-1/08-the-shell/outputs/run_check.txt}{Four good tables; a broken one and a missing one; no
argument.}\label{out:b1:the-shell:run-check}

\begin{method}[Writing a script that fails loudly]\label{met:b1:the-shell:script}
\begin{enumerate}
\item Start with \code{set -euo pipefail}: stop at the first failing command (\code{-e}), treat an unset variable as an
error (\code{-u}), and give a \omterm{def:b1:the-shell:pipeline}{pipeline} the status of the last of its commands that failed, or 0 if none did
(\code{pipefail}).
\item Check the arguments first; print the usage on standard error and exit with 2.
\item Quote every expansion: \code{"$f"}, \code{"$@"}, \code{"$(command)"}. An unquoted variable holding a name with a
space becomes two arguments.
\item Return a meaningful status: 0 when all is well, 1 for a negative answer, 2 for a misuse.
\item Where a failing status is expected, say so: \code{grep -c} exits 1 when it counts zero lines, which \code{-e}
would take for an error; \code{|| true} accepts it.
\end{enumerate}
\end{method}

\code{$(command)} is replaced by the output of the command, \code{$((...))} computes in integers, and
\texttt{\$\{\#commas\}} is the length of a variable's value. \code{if}, \code{for} and \code{while} test exit statuses:
\code{[ "$rows" -eq 0 ]} is the program \code{[}, which compares and exits with 0 or 1.

\section{Permissions and links}\label{sec:b1:the-shell:permissions}

\begin{definition}[Permission bits]\label{def:b1:the-shell:permission-bits}
The \emph{permission bits}\index{permission bits} of a file are nine bits, read (\code{r}), write (\code{w}) and execute
(\code{x}) for its owner, for the members of its group, and for all other users, written \texttt{rwxr-{}x-{}-{}-{}} by
\code{ls -l} or as three octal digits, one per class, each the sum of 4 (read), 2 (write) and 1 (execute): \code{750}.
For a directory, \code{r} allows listing it, \code{w} creating and removing names in it, and \code{x} entering it.
\end{definition}

\begin{omfigure}
\begin{tikzpicture}[x=6.5mm, y=1mm]
\newcommand{\permbit}[4]{\draw[fill=#3] (#1,0) rectangle (#1+1,7); \node[font=\small\ttfamily] at (#1+0.5,3.5) {#2};
  \node[font=\scriptsize] at (#1+0.5,-3) {#4};}
\permbit{0}{-}{black!8}{type}
\permbit{1}{r}{omCode!10}{4} \permbit{2}{w}{omCode!10}{2} \permbit{3}{x}{omCode!10}{1}
\permbit{4}{r}{omDef!12}{4} \permbit{5}{-}{omDef!12}{2} \permbit{6}{x}{omDef!12}{1}
\permbit{7}{-}{omThm!12}{4} \permbit{8}{-}{omThm!12}{2} \permbit{9}{-}{omThm!12}{1}
\foreach \a/\b/\t/\v in {1/4/owner/7, 4/7/group/5, 7/10/others/0} {
  \draw[decorate, decoration={brace, amplitude=3pt}] (\a,9) -- (\b,9) node[midway, above=3pt, font=\scriptsize]
  {\t: \v};}
\end{tikzpicture}
\par\smallskip{\small The mode \texttt{-{}rwxr-{}x-{}-{}-{}}, octal 750: the owner may read, write and execute; the group may read
and execute; others have no access. The first character is the type: \code{-} for a file, \code{d} for a directory,
\code{l} for a \omterm{def:b1:the-shell:links}{symbolic link}.}
\end{omfigure}

\omoutput{code/bachelor-1/08-the-shell/outputs/perms.txt}{\code{chmod} in symbolic and octal form, a refused write, and
the \code{umask}.}\label{out:b1:the-shell:perms}

\code{chmod go-r} removes read for group and others; \code{chmod 640} sets all nine bits at once. Even the owner cannot
write a file without its \code{w} bit, though the owner may give it back. A new file gets mode 666, a new directory
777, minus the bits of the \code{umask}: with 022, files are created \texttt{rw-{}r-{}-{}r-{}-{}}. Plugging in a removable drive
makes its file system appear as a directory of the tree, an operation called mounting, which requires the rights of the
administrator; desktop systems do it on the user's behalf.

\begin{definition}[Hard link, symbolic link]\label{def:b1:the-shell:links}
A directory entry is a name for a file. A \emph{hard link}\index{hard link}, made by \code{ln a b}, is a second name for
the same file: both names are equal, and the file disappears only when its last name is removed. A \emph{symbolic
link}\index{symbolic link}, made by \code{ln -s a c}, is a small file that holds a path; opening it opens the file at
that path, if there is one.
\end{definition}

\omoutput{code/bachelor-1/08-the-shell/outputs/links.txt}{The number after the mode is the number of names of the file;
the \omterm{def:b1:the-shell:links}{symbolic link} survives its target and leads nowhere.}\label{out:b1:the-shell:links}

\begin{omfigure}
\begin{tikzpicture}[e/.style={draw, font=\scriptsize\ttfamily, minimum width=22mm, minimum height=6mm},
                    d/.style={draw, rounded corners=3pt, fill=omDef!10, font=\scriptsize, align=center, minimum
                    width=24mm, minimum height=12mm}, >={Stealth[length=2mm]}]
\node[font=\scriptsize\bfseries] at (0,1.5) {directory entries};
\node[e] (s) at (0,0.8) {station.txt};
\node[e] (c) at (0,0.2) {copy.txt};
\node[e] (l) at (0,-0.4) {shortcut.txt};
\node[d] (f) at (4,0.5) {the file's data\\\code{Hilltop, 212 m}\\2 names};
\node[draw, dashed, font=\scriptsize, align=center] (p) at (4,-1.0) {a tiny file holding\\the path \code{station.txt}};
\draw[->, thick] (s.east) -- (f.west |- s.east);
\draw[->, thick] (c.east) -- (f.west |- c.east);
\draw[->, thick] (l.east) -- (p.west);
\draw[->, dashed, omThm] (p.north) -- node[right, font=\scriptsize] {by name} (f.south);
\end{tikzpicture}
\par\smallskip{\small Two \omterm{def:b1:the-shell:links}{hard links} and a \omterm{def:b1:the-shell:links}{symbolic link}. Removing \code{station.txt} leaves the file with one name,
\code{copy.txt}; the \omterm{def:b1:the-shell:links}{symbolic link} still holds the path \code{station.txt}, which now names nothing.}
\end{omfigure}

\begin{safety}
Quote every variable in a script. Test a destructive command with \code{echo} in front of it before running it, and
never write \code{rm -r} with a variable that might be empty. Keep your files \texttt{rw-{}r-{}-{}r-{}-{}} or stricter, and never
make a file writable by others to make an error go away.
\end{safety}

\section{Exercises}\label{sec:b1:the-shell:exercises}

\begin{exercise}[$\star$]\label{exo:b1:the-shell:1}
Predict what each line prints, with \code{weather.csv} the town's weather file.
\omcode{code/bachelor-1/08-the-shell/listings/exo_status.sh}{3}{6}{Four exit statuses}
\end{exercise}

\begin{exercise}[$\star$]\label{exo:b1:the-shell:2}
A directory holds \code{r1.csv}, \code{r2.csv}, \code{r10.csv}, \code{R3.csv}, \code{rain.txt} and \code{.r4.csv}. What do
\code{echo r?.csv}, \code{echo r*.csv}, \code{echo *.csv} and \code{echo [Rr][0-9]*} print?
\end{exercise}

\begin{exercise}[$\star$]\label{exo:b1:the-shell:3}
Predict the output of the three \omterm{def:b1:the-shell:pipeline}{pipelines}.
\omcode{code/bachelor-1/08-the-shell/listings/exo_pipe.sh}{3}{6}{Three \omterm{def:b1:the-shell:pipeline}{pipelines} on six lines}
\end{exercise}

\begin{exercise}[$\star$]\label{exo:b1:the-shell:4}
What does \code{ls -l} show as the mode of a file after \code{chmod 640}, \code{chmod 755}, \code{chmod 604}, and after
\code{chmod 640} followed by \code{chmod u+x,o+r}?
\end{exercise}

\begin{exercise}[$\star\star$]\label{exo:b1:the-shell:5}
The script below works until someone saves a file named \code{members 2024.csv}. Explain the run, and correct the
script.
\omcode{code/bachelor-1/08-the-shell/listings/count_lines.sh}{1}{5}{Counting lines}
\omoutput[2-5]{code/bachelor-1/08-the-shell/outputs/exo_space.txt}{The run.}
\end{exercise}

\begin{exercise}[$\star\star$]\label{exo:b1:the-shell:6}
Compare \code{ls nosuchfile > out1.txt 2>&1} with \code{ls nosuchfile 2>&1 > out2.txt}: where does the error message
go in each case, and why?
\end{exercise}

\begin{exercise}[$\star\star$]\label{exo:b1:the-shell:7}
Write a script \code{topk.sh FILE COL K} that prints the \code{K} most frequent values of column \code{COL} of a CSV file
with a header, each with its count, and rejects bad arguments with status 2.
\end{exercise}

\begin{exercise}[$\star\star$]\label{exo:b1:the-shell:8}
Write one \omterm{def:b1:the-shell:pipeline}{pipeline} that prints, for each station of the weather file, the line of its hottest day, ignoring missing
values and error codes.
\end{exercise}

\begin{exercise}[$\star\star$]\label{exo:b1:the-shell:9}
Among the 512 modes from \code{000} to \code{777}, how many let the owner read and let neither the group nor the others
write? Justify by counting.
\end{exercise}

\begin{exercise}[$\star\star\star$]\label{exo:b1:the-shell:10}
Write a script \code{dupes.sh DIR} that prints the groups of files with identical contents under \code{DIR}, using
\code{find}, the checksum program \code{cksum} (which prints a checksum and a size) and \omterm{def:b1:the-shell:filter}{filters}. It must handle names
with spaces. Why is a matching checksum not a proof of equal contents, and what would make the script faster on a
large tree?
\end{exercise}

\begin{exercise}[$\star\star\star$]\label{exo:b1:the-shell:11}
In an empty directory: \code{echo data > a; ln a b; ln -s a c; rm a}. Give the number of names of the file at each step,
and what \code{cat b} and \code{cat c} print at the end. Then \code{echo new > a}: what do they print now? Justify each
answer from the definitions.
\end{exercise}

\begin{exercise}[$\star\star\star$]\label{exo:b1:the-shell:12}
With \code{f="a b"}, \code{g='*'}, a directory holding only \code{x.txt}, and \code{empty} unset, how many arguments
does \code{printf} receive in each line below, and what does each line print? Justify with the order of expansions.
\omcode{code/bachelor-1/08-the-shell/listings/exo_words.sh}{4}{7}{Counting words}
\end{exercise}

\section{Problem: A day's log in one line}\label{sec:b1:the-shell:problem}

\begin{problem}[{Weekend problem --- A day's log in one line}]\label{pb:b1:the-shell:1}
The library's server writes a line to its access log for each request: the date and time, the member (or \code{-}
for an anonymous search), the action, the book, and a status, with a reason when it is an error. A generator with a
fixed seed reproduces one day of that log. The session below explores it with \omterm{def:b1:the-shell:filter}{filters}, answers two questions with a
\omterm{def:b1:the-shell:pipeline}{pipeline}, checks them with a Python program, and tries a script.
\omcode{code/bachelor-1/08-the-shell/listings/make_log.py}{3}{35}{The log generator}\label{lst:b1:the-shell:make-log}
\omcode{code/bachelor-1/08-the-shell/listings/logstat.sh}{1}{21}{A script for one hour of the log}\label{lst:b1:the-shell:logstat}
\omcode{code/bachelor-1/08-the-shell/listings/log_check.py}{3}{20}{The Python reference}\label{lst:b1:the-shell:log-check}
\omoutput{code/bachelor-1/08-the-shell/outputs/problem.txt}{The session.}\label{out:b1:the-shell:problem}

\needspace{4\baselineskip}
\medskip\noindent\textbf{Part I --- Filters and counts.}
\begin{enumerate}
\item How many requests did the library receive, and how many of them were searches?
\item Why does \code{cut -d" " -f4} give the action? What would \code{-f5} give?
\item What fraction of the requests failed?
\item Why does the \omterm{def:b1:the-shell:pipeline}{pipeline} of the actions not need \code{sort -rn} to be read, and why does it need \code{sort}?
\item Write the \omterm{def:b1:the-shell:pipeline}{pipeline} that counts the requests of member 24.
\item How many requests came between 12:00:00 and 13:59:59? Explain the pattern \code{"^.{11}1[23]:"}.
\end{enumerate}

\needspace{4\baselineskip}
\medskip\noindent\textbf{Part II --- Patterns.}
\begin{enumerate}[resume]
\item Explain the \code{sed} command of the error reasons, and why it is preceded by \code{grep status=error}.
\item The three reasons sum to how many lines? Check against question 3.
\item \code{grep -oE "member=[0-9]+"} prints only the matching part of each line. Why does it skip anonymous searches?
\item How many requests were anonymous, and how many were made by members? Check the total.
\item Give a \code{grep -E} pattern for the lines of an error during a borrow, and one for times from 18:30:00 on.
\item The \code{cut -c12-13} of the busiest-hour \omterm{def:b1:the-shell:pipeline}{pipeline} counts on fixed columns. Which line of the generator guarantees
that the hour is always at columns 12 and 13?
\end{enumerate}

\needspace{4\baselineskip}
\medskip\noindent\textbf{Part III --- A script with arguments and exit codes.}
\begin{enumerate}[resume]
\item What does \code{logstat.sh} check before reading the log, and with which status does it refuse?
\item Why does it distinguish \code{07} (status 1) from \code{7am} (status 2)?
\item The line \code{lines=$(grep -c ... || true)} is needed because of \code{set -e}. Explain.
\item \code{"^.{11}$hour:"} is in double quotes. What would change with single quotes?
\item The script reads the log twice. Rewrite it to read it once, with a temporary file or a single \omterm{def:b1:the-shell:pipeline}{pipeline}.
\item The chapter's test runs the script on the generated log for the busiest hour, for \code{07}, for \code{7} and with
one argument. Which statuses does it expect, and why is each case worth a test?
\end{enumerate}

\needspace{4\baselineskip}
\medskip\noindent\textbf{Part IV --- The pipeline against Python.}
\begin{enumerate}[resume]
\item Explain each stage of \code{cut -c12-13 access.log | sort | uniq -c | sort -rn | head -3}.
\item If two hours tied for the most requests, which one would \code{head -1} print, and which one would
\code{log_check.py}? Are they the same?
\item Why does the distinct-member \omterm{def:b1:the-shell:pipeline}{pipeline} need \code{sort -u} rather than \code{uniq}?
\item The generator chooses 42 visitors for the day. Why is it no surprise that 42 distinct members appear, and how
large would \code{n} have to be for one of them to be likely missing?
\item The \omterm{def:b1:the-shell:shell}{shell} and Python agree on this log. What does that agreement prove, and what does it not?
\item State the named result of the problem, \emph{the \omterm{def:b1:the-shell:pipeline}{pipeline} theorem of the log}.
\end{enumerate}
\end{problem}
